\documentclass[11pt,a4paper]{article}
\usepackage{freshaudit}
\fancyhead[L]{Critical audit --- Parts I and II}
\title{Fresh deep critical audit of Parts I and II\\[3pt]
\large Proof-breaking issues, repairs, and source-scope verification}
\author{Independent reconstruction pass}
\date{16 August 2026}

\begin{document}
\maketitle

\begin{resultbox}
\textbf{Bottom line.}

\textbf{Part I.}  The principal implication chain
\[
 \begin{aligned}
 \mathrm{GJS\ CE1/CE2}&\Longrightarrow M_2\Longrightarrow
 \text{uniform cutoff gap}\Longrightarrow\text{LPPL}\\
 &\Longrightarrow\text{full local-norm net limit}
 \Longrightarrow\text{spectral ground state}.
 \end{aligned}
\]
is mathematically coherent.  I found one omitted localization decomposition in the
CE2-to-gap step and one missing mollifier hypothesis in Appendix C.  Both have short,
fully explicit repairs.  The remaining substantive caveat is source scope: GJS
Theorem 1.1.8 is printed as uniform ``as $h\uparrow1$'', while Part I quantifies over
its entire enlarged cutoff class.  GJS p.~629 says the named estimates are uniform in
the space cutoff, which strongly supports the intended use, but the manuscript's
Appendix C does not reproduce enough of the cluster expansion to constitute an
independent proof of every strengthened quantifier.  The Main Theorem is therefore
\Conditional{} on that exact uniformity reading unless its cutoff class is narrowed.

\textbf{Part II.}  The current final theorem is \Fail{} as a proof in the files as
written.  The load-bearing shifted-interaction comparison (V2, Proposition 6.5) uses
an ambiguous cross power symbol, introduces an undefined $R_{\rm tail}$, and then
drops it.  This breaks A6 and hence gate (II.11).  Separately, V2 Proposition 1.2
contains a false combinatorial assertion: four edge momenta can conserve momentum.
That second error does \emph{not} change the advertised rate; its missing contribution
is explicitly bounded below.  The companion Lamport reconstruction supplies the
missing Fourier partition and a complete repair of Proposition 6.5, so the theorem is
salvageable, but the repair is new mathematics and must be adopted before Part II can
be called proved.
\end{resultbox}

\tableofcontents

\section{Scope, freshness, and standards}

\subsection{Documents of record}

The Part I document of record is \texttt{part\_i/main.tex}.  The current Part II proof
chain is the model sheet, free fixed-coarse rate, quadratic obstruction, projective
compactness, Weyl reduction, gate scoping note, and V1--V4.  The older M0--M6 files and
earlier audits were used only to locate possible source pages; no earlier verdict was
inherited.  Every cited page was rendered again from the local PDF.

\subsection{Validity standard}

A step passes only if its domains, quantifiers, normalizations, topology, and limiting
order are specified and every displayed equality follows from stated premises.  A true
result with a dropped term is still a proof defect.  A source citation passes only if
the exact cited source page states the needed result at the needed scope.  Numerical
checks are permitted only to falsify an identity, never to establish a theorem.

\subsection{Reproducibility artifacts}

The claim, dependency, source, citation, and finding ledgers are CSV files in
\texttt{ledgers/}.  The source-evidence companion contains 88 freshly rendered pages.
The Lamport reconstruction gives the repaired proof chain with hierarchical steps.

\section{Severity-ranked findings}

\subsection{F-II-02: the shifted-interaction comparison is not proved}

\begin{finding}[Dropped Fourier sectors in V2 Proposition 6.5]
\textbf{Severity: \Critical.}

V2 lines 1405--1417 assert a decomposition of the lattice, mixed, and continuum
quadratic forms, add a term called $R_{\rm tail}$, describe it only as ``regrouping''
missing modes, and then omit it from the inequality at lines 1419--1422.  In addition,
Lemma 6.4 defines the cross power symbol on the $j$-fold symmetric band but does not
put the edge weights $\chi$ into the defining convolution.  Its proof treats the
lattice symbol but does not separately prove the symmetric-edge cross case.

The omission is load-bearing: Proposition 6.5 is used by V4 Lemma 5.2, Theorem A6,
gate (II.11), and the final fixed-coarse theorem.
\end{finding}

\subsubsection{Why the displayed decomposition is not an identity}

Put $R=\pi/\eps_N$ and let
\[
 B_N=\{-R,-R+\pi/L,\ldots,R-\pi/L\},\qquad
 [k]_N\in B_N,\quad k-[k]_N\in2R\mathbb Z.
\]
Write $b_N(p)$ for the spatial Fourier coefficient of a lattice coefficient
$g_N$, $p\in B_N$, and $b_\infty(k)$ for the continuum coefficient.  If
$K_{\times,j}(\kappa)$ is the exact-output $j$-fold symbol of the translation-invariant
surrogate mixed kernel, direct Fourier expansion gives
\[
 \langle g_N,C_\times^jg_\infty\rangle
 =\frac1{2Lt}\sum_{\kappa}
 K_{\times,j}(\kappa)\,
 \overline{\widetilde b_N(-[k]_N,-k_0)}\,
 \widetilde b_\infty(k,k_0).
 \tag{2.1}
\]
The sum runs over the entire $j$-fold output band, not only $k\in B_N$.  Thus outputs
$k=p+2Rn$ with $n\ne0$ pair the lattice coefficient $b_N(-p)$ with the continuum
coefficient $b_\infty(p+2Rn)$.  They do not occur in the band-limited quadratic form
used at V2 line 1408.  Those are precisely mixed alias-output terms; calling them
$R_{\rm tail}$ does not make them vanish.

There is a second, independent correction.  The true mixed covariance is
$C_\times+R_N$ and is not translation invariant in $(x,y)$.  It may be bounded in
physical space, as V2 lines 1362--1388 do, but it cannot simultaneously be represented
by the unweighted cross symbol of Lemma 6.4.  One must choose one of the two descriptions
and account for the difference exactly once.

\subsubsection{A complete repair}

The following partition closes the gap.  It is proved step by step in the Lamport
reconstruction; the estimates are recorded here so the audit verdict is falsifiable.

Define, for $p\in (2\pi/t)\mathbb Z\times B_N$,
\begin{align*}
 K_{N,j}^{\rm fold}(p)
 &=(2Lt)^{-(j-1)}
 \sum_{\substack{\kappa_i\in D_N\\
 \sum k_{0,i}=p_0,\ [\sum k_i]_N=p}}
 \prod_i\widehat c_N(\kappa_i),\\
 K_{\times,j}(\kappa)
 &=(2Lt)^{-(j-1)}
 \sum_{\substack{\kappa_i\in D_N\\\sum\kappa_i=\kappa}}
 \prod_i\widehat c_\times(\kappa_i),\\
 K_{\infty,j}(\kappa)
 &=(2Lt)^{-(j-1)}
 \sum_{\sum\kappa_i=\kappa}\prod_i\widehat c_\infty(\kappa_i).
\end{align*}
All sums are nonnegative.  The envelope-convolution lemma implies, for
$j=1,2,3$,
\begin{align}
 \sup_p|K_{N,j}^{\rm fold}(p)-K_{\infty,j}(p)|
 &\le C\eps_N^2\log_N^2,\tag{2.2}\\
 \sup_{p\in B_N}|K_{\times,j}(p)-K_{\infty,j}(p)|
 &\le C\eps_N^2\log_N^2,\tag{2.3}\\
 \sup_{|k|\ge R}K_{\times,j}(k_0,k)
 &\le C\eps_N^2\log_N^{j-1}.\tag{2.4}
\end{align}
For (2.2), partition the folded lattice configurations into exact output and the
finitely many shifts $2Rn$.  Telescope the exact-output products one factor at a time,
use
\[
 |\widehat c_N-\widehat c_\infty|(u)
 \le C\eps_N^2\frac{u_{\rm sp}^4}{(m^2+|u|^2)^2},
\]
and sum the remaining $j-1$ envelope factors.  Every nonzero shift has one convolution
argument at distance at least $R$ and hence contributes (2.4).  Continuum configurations
with an input outside $B_N$ are bounded by fixing that input and applying the same tail
convolution.  This proves (2.2).  The proof of (2.3) is identical without folded output.
Equation (2.4) is the convolution-decay estimate at distance $R$.

Extend $b_N$ by zero outside the fundamental band and put
$\delta b=b_N-b_\infty$ on $B_N$, $\delta b=-b_\infty$ outside.  Splitting (2.1) into
$n=0$ and $n\ne0$ gives the exact identity
\begin{align*}
 Q_{N,j}-2Q_{\times,j}+Q_{\infty,j}
 &=\langle\delta b,K_{\infty,j}\delta b\rangle
 +\langle b_N,(K_{N,j}^{\rm fold}-K_{\infty,j})b_N\rangle\\
 &\quad-2\langle b_N,(K_{\times,j}-K_{\infty,j})b_{\infty,B}\rangle
 -2A_{\times,j},\tag{2.5}
\end{align*}
where
\[
 A_{\times,j}=\frac1{2Lt}
 \sum_{\substack{p_0\in(2\pi/t)\mathbb Z,\ p\in B_N,\ n\ne0}}
 K_{\times,j}(p_0,p+2Rn)\,
 \overline{b_N(-p,-p_0)}\,b_\infty(p+2Rn,p_0).
\]
Only finitely many $n$ occur because $K_{\times,j}$ is supported in the $j$-fold band.
Cauchy--Schwarz and (2.4) give
\[
 |A_{\times,j}|\le C\eps_N^2\log_N^2
 \|b_N\|_2\|b_\infty\|_2.\tag{2.6}
\]
The physical-space edge remainder satisfies
\[
 |Q_{\rm true\ cross}-Q_{\times,j}|
 \le C(2Lt)^2\Sigma_M(p)^{2(4-j)}\eps_N\log_N^{j-1}.
 \tag{2.7}
\]
Because $\eps_N=\eps_M2^{-r}$ and
$\theta_H=(1+\eta)/2<1$,
\[
 \eps_N\log_N^2\le C(\eps_M,\eta)2^{-r\theta_H}.
 \tag{2.8}
\]
Finally, the power coefficients obey
\[
 \|b_N^{(h_N^{4-j})}-b_\infty^{(h_\infty^{4-j})}\|_2^2
 \le C\,2^{-r\theta_H}.\tag{2.9}
\]
To prove (2.9), telescope the $(4-j)$-fold convolution; apply
$\ell^2*\ell^1\to\ell^2$ Young inequality to the ordinary differences; for lattice
aliases, $|k_1+\cdots+k_{4-j}|\ge R$ forces one
$|k_i|\ge R/(4-j)$, so put that factor in the $D4$ $\ell^2$ tail and the remaining
factors in $\ell^1$.  The same argument treats missing continuum modes.

Equations (2.2)--(2.9), inserted into the Wick--Isserlis identity, yield
\[
 \|\Delta\mathcal U_N-\Delta\mathcal U\|_2
 \le C\{2^{-r(1+\eta)/4}+\eps_N\log_N\},
\]
which is exactly the claimed rate.  Thus the proposition is repairable, but the repair
is absent from V2 and cannot be inferred from its undefined $R_{\rm tail}$.

\subsection{F-II-01: the ``three or more edge momenta are impossible'' claim is false}

\begin{finding}[Omitted four-edge sector]
\textbf{Severity: \Major.}

V2 Proposition 1.2(iii)(b) says that a fourth-moment cross configuration with at
least three edge momenta is impossible.  Let $R=\pi/\eps_N$.  The symmetric band contains
both $+R$ and $-R$, and
\[
 (+R)+(+R)+(-R)+(-R)=0.
\]
Hence six ordered four-edge sign configurations exist.  The exact cross-moment formula
at V2 lines 577--579 is correct; the subsequent classification is not exhaustive.
\end{finding}

Put $f(k_0)=\widehat c_\times(k_0,R)$; evenness makes the sign irrelevant.  Each edge
has weight $2^{-1/2}$, so the omitted positive contribution to
$\mathbb E[X_4Y_4]$ is
\begin{equation}
 \Delta_{4e}=4!(2Lt)^{-2}\frac32
 \sum_{k_{0,1}+\cdots+k_{0,4}=0}\prod_{i=1}^4f(k_{0,i}).
 \tag{2.10}
\end{equation}
The factor $3/2$ is
$\binom42(2^{-1/2})^4=6/4$.

Let $A=\sum f$ and $B=\sum f^2$.  Since
$f(k_0)\le(k_0^2+\gamma_e^2)^{-1}$,
$\gamma_e\ge2\eps_N^{-1}$, and $t\gamma_e/2\ge2t_0$,
\begin{align*}
 A&\le\frac{t}{2\gamma_e}\coth(t\gamma_e/2)
 \le\frac{t\eps_N}{4}\coth(2t_0),\\
 B&\le\gamma_e^{-2}A
 \le\frac{t\eps_N^3}{16}\coth(2t_0).
\end{align*}
Furthermore
\[
 \sum_{\sum k_{0,i}=0}\prod_i f(k_{0,i})
 =\|f*f\|_2^2
 \le\|f*f\|_\infty\|f*f\|_1
 \le B A^2.
\]
Substitution into (2.10) gives the explicit bound
\begin{equation}
 0\le\Delta_{4e}\le
 \frac{9t}{256L^2}\coth(2t_0)^3\eps_N^5.
 \tag{2.11}
\end{equation}
This is smaller than the $O(t\eps_N^3\log_N^2)$ correction already allowed.  Because it
increases the cross moment, dropping it also decreases the true squared difference
relative to the approximate expression.  Thus (2.11) repairs the equality and preserves
the rate; it does not by itself destroy A4(a).

\subsection{F-I-02: CE2 is applied to a nonlocalized polynomial}

\begin{finding}[Missing source-required cell decomposition]
\textbf{Severity: \Major in the written proof; short repair.}

Part I Lemma \texttt{ce2strip} applies GJS Theorem 1.1.7 directly to $Q$.  The source
theorem is stated for two localized monomials.  GJS p.~594 explicitly says that a
nonlocalized/general polynomial is written as a sum of localized monomials and defines
its norm as the sum of their norms.  Part I established this norm but omitted the
decomposition at the point of use.
\end{finding}

Write the finite source-prescribed decomposition
$Q=\sum_{a\in A}Q_a$.  Translation to either time row does not change the component
norms.  Every localization square of $Q_a(-t/2)$ is separated from every square of
$Q_b(t/2)$ by width at least $t-2$.  Therefore CE2 gives
\begin{align*}
 &\left|\operatorname{Cov}_{dq_h}
 \bigl(Q(-t/2),Q(t/2)\bigr)\right|\\
 &\quad\le\sum_{a,b}\left|\operatorname{Cov}_{dq_h}
 \bigl(Q_a(-t/2),Q_b(t/2)\bigr)\right|\\
 &\quad\le C_7e^{-m_1(t-2)}
 \sum_{a,b}\|Q_a\|_{\rm GJS}\|Q_b\|_{\rm GJS}\\
 &\quad=C_7e^{2m_1}e^{-m_1t}
 \left(\sum_a\|Q_a\|_{\rm GJS}\right)^2
 =C_7e^{2m_1}e^{-m_1t}\|Q\|_{\rm GJS}^2.
\end{align*}
This is exactly the displayed Lemma 7.4 estimate.  After insertion, semigroup decay and
the uniform-gap theorem are restored without changing a constant.

\subsection{F-II-03: V3's periodic N2 rate is asserted, not derived}

\begin{finding}[Unwritten two-dimensional tail estimate]
\textbf{Severity: \Major in a no-compression proof; repairable.}

V3 Proposition 2.6(iii) states that the periodic spatial truncations satisfy the N2 rate
by ``the same envelope tail counts,'' without writing the $k_0$ sum or the constrained
four-fold convolution.  V3 Lemma 6.1 then consumes that rate in the continuum
marked-slice limit.
\end{finding}

For the quartic chaos, Wick--Isserlis and a union bound over the first out-of-band input
give
\begin{align*}
 \|Y_4^{(K')}-Y_4^{(K)}\|_2^2
 &\le C(2Lt)^{-2}\sum_{|k|>K,k_0}
 \widehat c(k_0,k)F_3(-k_0,-k)\\
 &\le C\sum_{|k|>K,k_0}
 \frac{(1+\log(2+|k_0|+|k|))^2}
 {(m^2+k_0^2+k^2)^2}.
\end{align*}
For $t\in[t_0,T]$, the one-dimensional grid comparison gives
\[
 \sum_{k_0\in(2\pi/t)\mathbb Z}
 \frac{(1+\log(2+|k_0|+|k|))^2}{(m^2+k_0^2+k^2)^2}
 \le Ct\frac{(1+\log(2+|k|))^2}{(m^2+k^2)^{3/2}}.
\]
The spatial tail is bounded by
\[
 \sum_{|k|>K}\frac{(1+\log(2+|k|))^2}{(m^2+k^2)^{3/2}}
 \le CL\frac{(1+\log(1+K/m))^2}{K^2}.
\]
Thus the quartic norm is $O(\sqrt{tL}\,K^{-1}\log K)$, stronger than V3's
$O(K^{-1/2}\log^2K)$.  The quadratic chaos uses
$\sum_{|k|>K,k_0}\widehat c(k_0,k)^2$ and has the same or better rate.  The difference
between the truncated thermal-to-vacuum Wick constants is exponentially small because
\[
 0\le\frac{coth(t\gamma/2)-1}{2\gamma}
 =\frac{1}{\gamma(e^{t\gamma}-1)}\le C(t_0,m)e^{-t_0\gamma}/\gamma.
\]
The claimed N2 estimate follows.  This derivation is included in the reconstruction.

\subsection{F-I-01: GJS cutoff uniformity is stronger than the theorem's literal wording}

\begin{finding}[Source-quantifier ambiguity]
\textbf{Severity: \Major if full-class uniformity is not accepted; otherwise citation
qualification.}

GJS Theorem 1.1.7 is explicitly independent of $h$.  Theorem 1.1.8 is printed as a
bound holding uniformly ``as $h\uparrow1$.''  Part I promotes it to every
$h\in\mathcal C_E$, including arbitrary compact slabs and affine cutoff paths.  GJS
p.~629 states that the estimates of Theorems 1.1.8, 1.1.11, 3.1 and 4.1 are uniform in
the space cutoff $h$.  This is strong positive support, but it is not a verbatim
restatement of the enlarged quantifier used in Part I.
\end{finding}

Appendix C's ``structural uniformity'' proposition lists the locations where $h$ enters
the cluster expansion, but it does not reproduce the polymer estimates, convergence
summations, or constants.  It is therefore a source audit, not a standalone proof of a
strengthened cluster-expansion theorem.  There are two rigorous ways to close the issue:

\begin{enumerate}
\item state as an explicit imported assumption that GJS p.~629's ``uniform in the space
cutoff $h$'' covers the entire source cutoff class used here; or
\item narrow the manuscript cutoff family to a directed family covered literally by the
printed limiting statement and prove that every slab/path application belongs to a
uniform tail of that family.
\end{enumerate}

The present manuscript does neither with complete formal precision.  This is why the
Part I verdict is conditional rather than an unconditional fail: the source evidence is
substantial and the internal argument needs only the bound, but the exact import scope
must be made explicit.

\subsection{F-I-03: the ultraviolet mollifier lacks used hypotheses}

\begin{finding}[Approximate identity underspecified]
\textbf{Severity: \textbf{MODERATE}.}

Part I Appendix C fixes a ``smooth approximate identity'' $\delta_\kappa$ and later
treats $G^{(1)}_\kappa$ and $G_\kappa$ as averages of $G$ over balls of radius at most
$2/\kappa$.  This requires $\delta\ge0$, $\int\delta=1$, compact support, the scaling
$\delta_\kappa(x)=\kappa^2\delta(\kappa x)$, and for the symmetric Wick-covariance
identities it is convenient to require $\delta(x)=\delta(-x)$.  None is stated.
\end{finding}

Choose once and for all
$\delta\in C_c^\infty(B(0,1))$, $\delta\ge0$, $\delta(x)=\delta(-x)$,
$\int\delta=1$, and define $\delta_\kappa(x)=\kappa^2\delta(\kappa x)$.  Then the
``average over balls'' bounds and all covariance identities used in the proof follow.
This repairs the lemma.  It does not independently repair the broader cluster-expansion
uniformity issue F-I-01.

\subsection{F-II-04: a constant has unstated anisotropic-grid dependence}

V2 Lemma 3.1 describes certain convolution constants as depending only on $m$.  The
proof absorbs isolated terms by $tL\ge1$ and uses a lower bound $t\ge t_0$; hence the
constant also depends on the grid control, equivalently on $t_0$ (and on fixed $L$ if
the condition is not normalized).  Downstream statements already allow that dependence,
so this is not proof-breaking.  Correct the displayed dependency.

\subsection{F-II-05: finite-mode transplant is too compressed for the stated standard}

V3 Theorem 3.1 transfers the model-sheet and V1 finite-dimensional kernel results into
the continuum-mode truncations.  The spaces are finite-dimensional, the covariance
matrices are positive, and the Mehler/Trotter argument is valid; no contradictory
normalization was found.  Nevertheless, the theorem repeatedly says that earlier steps
are ``transplanted'' instead of writing the unitary coordinate identification, reference
Gaussian density, and trace-kernel normalization.  The Lamport reconstruction supplies
those steps.  This is a rigor/presentation defect, not an identified false result.

\section{Part I proof-chain audit}

\subsection{Imports and dimensionless reduction}

The cutoff Hamiltonian existence/uniqueness, FKN formula, stability, hypercontractivity,
domain smoothing, finite propagation, and covariance claims have exact positive source
pages in the evidence dossier.  Three former pinpoint risks are resolved:

\begin{enumerate}
\item Simon III.6 and V.14--V.15 are the FKN results; III.12 is not.
\item Simon I.17 and III.17 are the hypercontractive results; V.10 is not the direct
source of the displayed abstract inequality.
\item Glimm--Jaffe Expositions Theorem 7.4 is the model-specific domain-smoothing theorem;
Theorem 7.3 alone has an extra bounded-below premise.
\end{enumerate}

The dilation in Part I tracks the mass, coupling, spatial cutoff, Wick covariance, and
observable norm consistently.  No reversed small-coupling inequality was found.

\subsection{Slab limit and uniform moment}

The bounded FKN formula is extended first to signed bounded functions, then to unbounded
polynomials by truncation and uniform square moments.  Simplicity alone identifies the
large-slab rank-one limit.  The $M_2$ proof uses CE1 on $V(h)^2$, which is permitted as a
product of individually Wick-ordered monomials; no Wick recombination is needed.  The
single overall $K_1$ in the GJS norm is used correctly.

\subsection{Uniform gap}

After inserting F-I-02's cell decomposition, the logic is sound.  The Euclidean two-time
bound gives
$\langle\psi_Q,e^{-tH(g)}\psi_Q\rangle\le C_Qe^{-m_1t}$ on a dense family.
If $P_{(0,m_1)}\ne0$, choose a vector in its range and approximate it by that dense
family; the spectral theorem gives a lower exponential rate incompatible with the
upper rate.  Simplicity identifies the zero-energy projection.  The proof does not infer
a gap from clustering for only one observable; it first proves density, which is the
necessary step.

\subsection{Path regularity, filter, and LPPL}

The affine path stays within the admissible cutoff class.  The semigroup derivative is
proved in a weak pairing and upgraded using the uniform $M_2$ control.  The spectral
projection derivative is licensed by the common uniform gap.  The exact filter has the
correct Fourier convention and equals $H^{-1}(1-P)$ on the gapped subspace.  Finite
propagation localizes the truncated filter, while the time tail is controlled by the
filter decay.  The unbounded interaction is used only through affiliation and the
uniform second moment.  Shell counting in one space dimension produces a summable,
super-polynomial LPPL rate.  No hidden assumption of bounded interaction was found.

\subsection{Full net, normality, and spectral passage}

The plateau-containment relation is technically a preorder rather than an antisymmetric
partial order, but directed-net convergence is unaffected.  Local dual spaces are
complete, restriction limits are compatible, and positivity extends to the quasi-local
completion.  Normality is proved from norm convergence to normal cutoff vector states;
the monotone-net argument is valid.  The ground-state spectral condition is passed using
exact equality with a nonnegative cutoff spectral measure on every fixed compact time
window, followed by an $L^1$ tail limit.  This avoids an unjustified convergence of
Hamiltonians.

\subsection{Dichotomy scope}

The two-phase exclusion theorem is explicitly conditional on hypothesis B4.  The cited
GJS phase papers prove phase coexistence and pure-phase gaps in their stated regimes but
do not supply the classification of every parity-invariant Hamiltonian ground state
required by B4.  Because the manuscript now labels B4 as a hypothesis, this is not a
false theorem; it must not be advertised as an unconditional uniqueness result.

\section{Part II proof-chain audit}

\subsection{Correct target and free obstruction}

The quadratic obstruction correctly rejects a diagonal finest-scale local-state norm
target.  MMST's construction is projective: convergence is evaluated after pullback to
one fixed coarse algebra.  The localized high-frequency sequence exposes the difference
between pointwise/fixed-coarse weak-star convergence and a supremum over a moving unit
ball.  The projective compactness theorem and full-sequence Weyl criterion are standard
compactness/density arguments and are correct.

The exact D4 finite product, dispersion comparison, and fixed-coarse free rate were
recomputed.  The support, normalization, and momentum/position covariance factors agree
with MMST.  The cited source proves convergence but not the new rate; the rate is supplied
internally and its summations are adequate.

\subsection{V1: thermal representations}

The finite-dimensional Schr\"odinger realization produces the correct Weyl phase.  The
Mehler kernel normalization integrates to the harmonic trace.  Completing the square in
the bridge variables yields the stated harmonic jump interpolant and action.  The V4
erratum correctly specifies the periodic Gaussian reference measure, avoiding a circular
use of the interacting measure.  Mode products converge under the D4 envelope, and all
uses of the thermal rate occur for $t\ge t_*>0$.

\subsection{V2: unshifted interaction and Wick dictionary}

The same-noise coupling needs the special edge symmetrization because the half-open
spatial band is not invariant under $k\mapsto-k$.  The marginal lattice covariance is
correct.  The unshifted A4(a) comparison has valid dispersion, missing-mode, and lattice
alias estimates once the four-edge term (2.10) is added.  The Wick recombination formula
and the vacuum/thermal dictionary constants have the correct signs; no residual mass
counterterm remains under the matched convention.

The shifted comparison is the unique unresolved load-bearing point in the original file;
see F-II-02.

\subsection{V3: continuum Hamiltonian package}

The finite-mode FKN semigroups are uniformly lower bounded by Nelson estimates.  Strong
semigroup convergence defines a self-adjoint generator; the cylinder-core generator
identity and lower form bound are compatible.  The hypercontractive ladder upgrades a
single-time Hilbert--Schmidt factor to trace class and permits trace-norm finite-mode
convergence.  Eigenvalue and partition-function convergence then follow from standard
trace-ideal arguments.  The periodic FKN clause requires the quantitative tail repair in
F-II-03.  Positivity improving follows from Simon I.16 because
$\|e^{-t\gamma}\|=e^{-tm}<1$; simplicity and a positive first gap follow from compact
resolvent/trace class.  No use at norm $1$ was found.

\subsection{V4: assembly and quantifiers}

Every invocation of V2's $tL\ge1$ estimates occurs with
$t\ge\max(1,1/L)$.  The corrected Laplace-transfer lemma needs transform convergence
only for $t\ge t_*$, and V4 uses exactly that range.  The interchange lemma requires a
uniform spectral tail, supplied by the transferred gap and normalized partition bound.
Thus the order
\[
 \text{fixed }t:\ N\to\infty,\qquad\text{then }t\to\infty
\]
is legitimate once A6 is proved.  Because A6 currently depends on the incomplete V2
Proposition 6.5, the final theorem is not proved in the original chain.  With the repaired
Proposition 6.5 and F-II-03 inserted, the remaining assembly closes.

\section{External-source verdicts}

\begin{longtable}{P{3.0cm}P{3.0cm}P{4.0cm}P{4.0cm}}
\toprule
Source/result & Exact evidence & Positive conclusion & Limitation\\
\midrule\endhead
GJ I--II & I02--I08 & cutoff core, ground state, covariance, locality & no full-net
vacuum convergence\\
GJ Expositions 7.4 & I24--I25 & model-specific semigroup domain smoothing & 7.3 alone
has extra premise\\
GJS 1.1.7 & I31 & CE2 and explicit $h$ independence & localized monomials; decompose
general $Q$\\
GJS 1.1.8/p.629 & I32--I33 & CE1 and space-cutoff uniformity mechanism & full-class
quantifier should be stated as import\\
GRS II.14/II.16 & I34--I35 & FKN and general marked-time FKN & none relevant\\
Simon I.17/III.17 & I44--I45/I48--I49 & abstract and time-region hypercontractivity &
V.10 is not the direct pinpoint\\
Simon V.7/V.14--15 & I50/I54--I55 & stability and interacting FKN & none relevant\\
MMST 4.1--4.5 & II07--II10 & observable embedding and free fixed-coarse limit & no
interacting theorem or rate\\
MMST Section 6.1 & II14 & identifies interacting extension as future work & cannot be
used as theorem\\
Battle--Federbush & II15--II17 & phase-cell wavelet mechanism & different variables;
no A4/A5 estimates\\
NRSS 4.1 & II19 & anharmonic LR under Fourier/$L^1$ assumptions & Wick quartic fails
$V'\in L^1$\\
\bottomrule
\end{longtable}

\section{Final disposition}

\begin{longtable}{P{3.1cm}P{3.4cm}P{7.5cm}}
\toprule
Deliverable theorem & Status & Exact disposition\\
\midrule\endhead
Part I Main Theorem & \Conditional & Internal chain closes after the explicit CE2
decomposition.  Accept GJS p.~629 as a full cutoff-class import or narrow the class.
State the mollifier hypotheses in Appendix C.\\
Part I canonical-vacuum corollary & \Conditional & Follows from the Main Theorem and
the verified GJ III construction formulas.\\
Part I two-phase exclusion & \Conditional & Correct only under explicit B4; cited
phase papers do not discharge B4.\\
Part II unshifted A4(a) & \Repair & Add the four-edge contribution (2.10)--(2.11).
The rate survives.\\
Part II shifted A4(c)/Proposition 6.5 & \Fail{} as written & Undefined/dropped mixed
alias sectors break the proof.  Equations (2.1)--(2.9) provide a complete replacement.\\
Part II A6 and gate (II.11) & \Fail{} as written & Depend on Proposition 6.5.  They
become valid after the replacement and the explicit V3 periodic-tail derivation.\\
Part II final fixed-coarse theorem & \Fail{} as written; salvageable & The free target,
compactness, Weyl reduction, V1, and V3/V4 architecture are sound.  Adopt the two
repairs before asserting the theorem.\\
\bottomrule
\end{longtable}

No evidence was found that the cited literature already proves the new interacting
fixed-coarse Part II comparison.  The value of the present audit is therefore not a
source substitution: it isolates the exact internal Fourier estimate on which the new
result stands or falls and supplies a checkable repair.

\end{document}
